\section{自我提升的推理路径}
\subsection{从 Chain-of-Thought 到自我奖赏}
Weston 强调，推理能力的增长不能仅靠强 prompt，而需要模型本身在训练阶段 ``思考自己的思考''。Self-Rewarding language models 由此而生：模型在训练时同时扮演题目生成者、解答者与评审，用自我评分来代替人类打分，就像把 Self-Feeding Chatbot 的 ``verified response'' 机制扩展到推理任务。这个过程要消除 ``记忆 vs 推理'' 的混淆，因此 base data 仍包含强制 chain-of-thought 的样本，以便 reward model 能识别出哪一段 reasoning 解释了答案。

\subsection{推理与自我评价的平衡}
模型在 self-improvement 迭代中容易陷入过度链式推理：它为了获得更高 reward 会生成越来越长但不一定更正确的 chain。Weston 引入 ``self-rewarding vs self-feeding'' 的对照，指出后者依赖对话反馈，而前者依赖自我验真（self-verification）。他建议在 self-rewarding 训练中插入 ``cross checking''：如果模型在初始 draft 中没能拒绝错误，就让它再 ``问自己'' 一次。

\begin{warningbox}{避免 overthinking}
出于安全考虑，必须限制 Chain-of-Thought 的膨胀，否则 inference compute 会被浪费在重复的思考上。Meta judge 的引入部分就是为了及时发现 ``即便有长推理链也答错'' 的情况。
\end{warningbox}

\subsection{训练信号与奖励模型}
在讲座中，Weston 用三类 reward 概念去解释信号来源：Process Reward Model（PRM，逐步奖励）、Outcome Reward Model（ORM，只看最终答案）以及传统的 RLHF / DPO（人类偏好）。Self-rewarding 训练更接近 PRM，它会在模型每一步推理之后尝试估分，而 Meta judge 则用 ORM 视角对比 draft 与 verified response。Weston 借用 ``DeepSeek-R1、O1-R1'' 这样的 benchmark 例子说明：在没有额外人工评审的情况下，自我奖励仍能把 reasoning win rate 推高（尤其在 Alpaca Eval 2 上）。

\begin{importantbox}{奖励模型对比}
\begin{itemize}
    \item PRM（过程奖励）：每一步 chain 都有信号，适合调教 reasoning trace。
    \item ORM（结果奖励）：只看最终判断，适合 meta judge 这类只需要判断哪个回答更好。
    \item RLHF / DPO：需要人工偏好，训练成本高，作为 baseline 仍被用于对照。
\end{itemize}
\end{importantbox}

\begin{knowledgebox}{可验证奖励的价值}
Weston 点出 ``verifiable reward'' 的理念：如果模型能把自己的 verified response 视为一个近似 ground truth，就可以使用这些 data 来 train reward 模型，强化 self-evaluation，而非完全依赖人工好坏标签。
\end{knowledgebox}

\subsection{本章小结}
自我提升的核心在于重写 reward signal：让模型自己判定推理链的价值，并在训练中同时保留 process/ outcome view。后续章节会把这一层的训练目标具体化成可复用的 pipeline。

